\emph{Conventions.} The modes of distinct channels and levels are distinct tower planes, $\{c_{\alpha}, c^{*}_{\beta}\} = \delta_{\alpha\beta}$, all other anticommutators zero (Fact \claimref{field-theory.imported-tower-limit}); the check realises them by Jordan-Wigner signs and verifies the Clifford relations of the planes $e = c^{*} + c$, $f = c^{*} - c$. $J$ commutes with every mode. The pair currents are those of Theorem \claimref{synthesis.null-plane-photon-data-and-the-pair-currents-of-the-generators}(iv), $j^{xy}_{-q} = \sum_{k}c^{*}_{x,k+q}c_{y,k}$ for $q \ge 1$.

\emph{(i)} For $q \ge 0$, moving $c_{y,k}$ past $c^{*}_{z,k'}$,
\[
j^{xy}_{-q}j^{zu}_{q} = \delta_{yz}\sum_{k}c^{*}_{x,k+q}c_{u,k+q} - \sum_{k,k'}c^{*}_{x,k+q}c^{*}_{z,k'}c_{y,k}c_{u,k'+q},
\]
sums over the levels in the window. With the weights $2$ for $q \ge 1$ and $1$ for $q = 0$, the first terms give $\sum_{x,u}(K_{C})_{xu}\sum_{l}(2l + 1)c^{*}_{x,l}c_{u,l}$, since a level $l$ is reached from $q = 1, \dots, l$ with weight $2$ and from $q = 0$ with weight $1$. Label a quartic term $c^{*}_{\alpha}c^{*}_{\beta}c_{\gamma}c_{\delta}$ by four modes: here $\alpha = (x, k + q)$, $\beta = (z, k')$, $\gamma = (y, k)$, $\delta = (u, k' + q)$, with coefficient $W(\alpha, \beta, \gamma, \delta) = C_{x_{\alpha}x_{\gamma}, x_{\beta}x_{\delta}}$, and each quadruple with $k_{\alpha} + k_{\beta} = k_{\gamma} + k_{\delta}$ and $k_{\alpha} \ge k_{\gamma}$ arises exactly once ($q = k_{\alpha} - k_{\gamma}$). So the quartic part of $\Phi(C)$ is $-(2X + D)$, $X$ the sum of $W\,c^{*}_{\alpha}c^{*}_{\beta}c_{\gamma}c_{\delta}$ over $k_{\alpha} > k_{\gamma}$ and $D$ over $k_{\alpha} = k_{\gamma}$. The relabelling $(\alpha, \beta, \gamma, \delta) \mapsto (\beta, \alpha, \delta, \gamma)$ maps the quadruples with $k_{\alpha} < k_{\gamma}$ (then $k_{\beta} > k_{\delta}$) onto those with $k_{\alpha} > k_{\gamma}$, leaves the product unchanged (two transpositions), and leaves $W$ unchanged because $C_{x_{\beta}x_{\delta}, x_{\alpha}x_{\gamma}} = C_{x_{\alpha}x_{\gamma}, x_{\beta}x_{\delta}}$. Hence $2X + D = S$, the sum of $W\,c^{*}_{\alpha}c^{*}_{\beta}c_{\gamma}c_{\delta}$ over all level-conserving quadruples. Exchanging the names $\gamma$ and $\delta$ preserves that index set and reverses $c_{\gamma}c_{\delta}$, so $S = \tfrac12\sum(W(\alpha, \beta, \gamma, \delta) - W(\alpha, \beta, \delta, \gamma))c^{*}_{\alpha}c^{*}_{\beta}c_{\gamma}c_{\delta} = \tfrac12 Q(A_{C})$, which gives (i); this is the antisymmetrisation argument of the proof of Theorem \claimref{synthesis.exact-window-sugawara-identity-and-luttinger-normal-modes-of-tower-density-waves}(i), with the channel tensor carried along. The check verifies (i) on every basis state for three generic symmetric tensors with $A_{C} \ne 0$ ($(n, L) = (2, 3), (2, 2), (3, 2)$). For $C = \sum_{A}M^{A}\otimes M_{A}$ the tensor is symmetric because the Casimir element is ($\sum_{A}M^{A}\otimes M_{A} = \sum_{A}M_{A}\otimes M^{A}$, the trace-dual of the trace-dual being the original basis), and $\Phi(C)$ is the stated bilinear in the currents by linearity.

\emph{(ii)} $C_{xy,zu} = \delta_{yz}\delta_{xu} + \delta_{xy}\delta_{zu}$ is symmetric, $(K_{C})_{xu} = \sum_{y}(\delta_{yy}\delta_{xu} + \delta_{xy}\delta_{yu}) = (n + 1)\delta_{xu}$, and $A_{C}(a, b, c, d) = \delta_{cb}\delta_{ad} + \delta_{ac}\delta_{bd} - \delta_{db}\delta_{ac} - \delta_{ad}\delta_{bc} = 0$. So $\Phi(C) = 2(n + 1)H$, and since $\Phi(C)$ is $2\sum_{q\ge1}(\sum_{x,y}j^{yx}_{-q}j^{xy}_{q} + N_{-q}N_{q}) + \sum_{x,y}j^{xy}_{0}j^{yx}_{0} + N_{0}^{2}$ (relabel $x \leftrightarrow y$ in the first term) the display follows. For $n = 1$ both bilinears are the stored $\rho_{q}^{*}\rho_{q}$ and $N_{0}^{2} = \hat N^{2}$, so $4H = 4T + 2\hat N^{2}$, that is $\sum_{j}(j + \tfrac12)n_{j} = T + \tfrac12\hat N^{2}$, the stored identity $T = \sum_{j}j\,n_{j} - \tfrac12\hat N(\hat N - 1)$. Since $\sum_{x,y}E_{xy}\otimes E_{yx} + 1\otimes 1 = C_{\su(n)} + \tfrac{n + 1}{n}1\otimes 1$ ($C_{\su(n)}$ the Casimir tensor of the trace-orthogonal complement of $1$) and $\sum_{x,y}[E_{yx}, [E_{xy}, X]] = 2nX - 2\operatorname{tr}(X)1$ gives $\lambda_{\su(n)} = 2n$, $H = \Phi(C_{\su(n)})/(2n + 2) + \Phi(1\otimes 1)/(2n) = H_{\su(n)} + H_{N}$ with $H_{N} = \tfrac12\Phi(\tfrac1n 1\otimes 1)$. The check verifies (ii) on every basis state for $(n, L) = (1, 4), (2, 3), (3, 3)$ and on three states of sixteen channels, and that the coefficient $n$ in place of $n + 1$ fails.

\emph{(iii)} The sixteen channels of a light ray and generation are the Weyl components valued in the particle half (Theorem \claimref{synthesis.slot-module-light-ray-arrays-are-weyl-components}(ii)). The pair bilinears change the mode number by $\pm 2$ and preserve the even half (Theorem \claimref{synthesis.spin10-clause-content-is-the-neutral-singlet-choice}); the diagonal generators $a^{*}_{a}a_{a} - \tfrac12$ have trace $8 - 8 = 0$ there, the others are off-diagonal. With the trace-dual basis computed exactly, the check finds: $\sum_{A}M^{A}M_{A} = \tfrac{45}{16}1$ (the trace of the left side is $\dim = 45$, and it is scalar), $A_{C} = -\tfrac{3}{16}\Delta$ with $\Delta(a,b,c,d) = \delta_{ac}\delta_{bd} - \delta_{ad}\delta_{bc}$ on all $16^{4}$ index quadruples, and $\lambda_{\so(10)} = 4$. (The reason $A_{C}$ is a multiple of $\Delta$: $A_{C}$ is the matrix of the operator $\sum_{A}M^{A}\otimes M_{A}$ on $16\otimes 16$ restricted to the antisymmetric part $\Lambda^{2}$, where it equals $\tfrac12$ of the Casimir of $\Lambda^{2}$ minus those of the two factors, and $\Lambda^{2}$ of the half-spinor of $\so(10)$ is irreducible, of dimension $120$; this is argued, the multiple is computed.) For $1\otimes 1/16$, $K = \tfrac1{16}1$ and $A = \tfrac{1}{16}\Delta$. By (i),
\[
H_{\so(10)} = \tfrac16\Bigl(\tfrac{45}{8}H + \tfrac{3}{32}Q(\Delta)\Bigr) = \tfrac{15}{16}H + \tfrac{1}{64}Q(\Delta), \qquad H_{N} = \tfrac12\Bigl(\tfrac18 H - \tfrac{1}{32}Q(\Delta)\Bigr) = \tfrac{1}{16}H - \tfrac{1}{64}Q(\Delta),
\]
using $\sum_{l}(2l + 1)n_{x,l} = 2\sum_{l}(l + \tfrac12)n_{x,l}$, and the sum is $H$. By (ii) with $n = 16$, $H - H_{N} = H_{\su(16)}$. The central values $2\dim/(2 + \lambda)$ are $90/6 = 15$ and $510/34 = 15$. The check also verifies $H = H_{\so(10)} + H_{N}$ directly, by applying the current bilinears, on three arbitrary states of a two-level window and on the sea, a state of charge $+1$, one of charge $-1$ and a particle-hole state of a six-level window.

\emph{(iv)} For the $25$ number-conserving generators alone the check finds that $A_{C}$ has entries off the support of $\Delta$, so it is not a multiple of $\Delta$ and (i) does not reduce their Sugawara energy to a multiple of $H$ plus a multiple of $Q(\Delta)$. For $X$ traceless in $\gl(5)$, $\sum X_{ab}a^{*}_{a}a_{b} = \sum X_{ab}(a^{*}_{a}a_{b} - \tfrac12\delta_{ab})$ lies in the span of the $45$ generators, so the gauge algebra of Theorem \claimref{synthesis.cpt-closure-forces-slot-module-and-chirality} lies in this $\so(10)$; the displayed equality is (iii) with $\sum_{i}H_{\mathfrak{f}_{i}}$ subtracted from both sides.
